\documentclass[a4paper,12pt]{article}
\usepackage{titlesec}
\titleformat*{\section}{\Large\bf}
\titleformat*{\subsection}{\normalsize\bf}
\titleformat*{\subsubsection}{\normalsize\bf}
%
%\usepackage[latin1]{inputenc}
\usepackage[utf8]{inputenc}
\usepackage[swedish,english]{babel}

%	Nationella tecken måste kodas
%	å=\aa,  ä=\"{a},   ö=\"{o} 
%	Å=\AA, Ä=\"{A}, Ö=\"{O} 


\addto\captionsenglish{%
  \renewcommand{\listoffigures}{Figurer}%
  \renewcommand{\listoftables}{Tabeller}%
  \renewcommand{\contentsname}{Inneh\aa llsf\"{o}rteckning}%
}

%\usepackage [T1] {fontenc}
%\usepackage [swedish] {babel}
\usepackage[numbers]{natbib}
\usepackage{float}
\usepackage{my.techs}
\usepackage{graphicx}
\usepackage{hyperref}

% \Section ger mindre spillutrymme, anv�nd dem om du vill
\newcommand{\Section}[1]{\section{#1}\vspace{-8pt}}
\newcommand{\Subsection}[1]{\vspace{-4pt}\subsection{#1}\vspace{-8pt}}
\newcommand{\Subsubsection}[1]{\vspace{-4pt}\subsubsection{#1}\vspace{-8pt}}

\hypersetup{
    %bookmarks=true,         % show bookmarks bar?
    %unicode=false,          % non-Latin characters in Acrobat�s bookmarks
    %pdftoolbar=true,        % show Acrobat�s toolbar?
    %pdfmenubar=true,        % show Acrobat�s menu?
    %pdffitwindow=false,     % window fit to page when opened
    %pdfstartview={FitH},    % fits the width of the page to the window
    pdftitle={Project title},    % title
    pdfsubject={Computer Science}, % subject
    pdfauthor={Name of author},     % author
    pdfkeywords={Here you write the keywords related to the theses}, % list of keywords
    pdfnewwindow=true,      % links in new window
    plainpages=false,				 % false: allows refs to pages that repeat (i.e. 1 in the preface as well as report)
    colorlinks=true,       % false: boxed links; true: colored links
    linkcolor=black,          % color of internal links
    citecolor=black,        % color of links to bibliography
    urlcolor=blue
    %filecolor=magenta,      % color of file links
}
\def\inst{datavetenskap}
\def\typeofdoc{Laborationsrapport}
\def\course{Examensarbete 15p}
\def\pretitle{Rapport under utveckling}
\def\title{Rapporttitel}
\def\name{Fnamn Enamn}
\def\username{anv\"{a}ndarnamn}
\def\email{\username{}@student.lnu.se}
\def\path{edu/KURS/lab1}
\def\graders{HANDLEDARNAS NAMN}


% om du vill referera till katalogen d�r dina filer ligger kan du 
% anv�nda \fullpath som kommer att vara "~username/edu..." o.s.v.
%\def\fullpath{\raisebox{1pt}{$\scriptstyle \sim$}\username/\path}


\begin{document}
	\pagenumbering{roman}
	%Here the preface - titlepage, TOC etc - starts:
	%---------------------------------------------------------------------------------------------------------------
%	\title{Project title}
\subtitle{- subtitle}
\author{Name of author}
\semester{Semester}
\subject{Computer Science}
\level{Bachelor/Master}
\coursecode{CourseCode}
\createtitlepage

	% skapar framsidan (om den inte duger: g�r helt enkelt en egen)
	\begin{titlepage}
		\thispagestyle{empty}
		\begin{large}
			\begin{tabular}{@{}p{\textwidth}@{}}
				\textbf{Linn\'{e}universitetet \hfill \today} \\
				\textbf{Institutionen f\"{o}r \inst} \\
				%\textbf{\typeofdoc} \\
			\end{tabular}
		\end{large}
		\vspace{10mm}
		\begin{center}
			\LARGE{\pretitle} \\
			\LARGE{\course}\\
			\vspace{10mm}
			\LARGE{\textbf{\title}} \\
			\vspace{15mm}
			\begin{large}
				\begin{tabular}{ll}
					\textbf{Namn} & \name \\
					\textbf{E-mail} & \texttt{\email} \\
					%\textbf{S�kv�g} & \texttt{\fullpath} \\
				\end{tabular}
			\end{large}
			\vfill
			\large{\textbf{Handledare}}\\
			\mbox{\large{\graders}}
		\end{center}
	\end{titlepage}



	%\title{Project title}
\subtitle{- subtitle}
\author{Name of author}
\semester{Semester}
\subject{Computer Science}
\level{Bachelor/Master}
\coursecode{CourseCode}
\createtitlepage
	%%%%%%%%%%%%%%%%%%%%%%%%%%%%%%%%%%%%%%%%%%%%%%%%%%%%%%%%%%%%%%%%%%%%%%%%%%%%%%%%%%%%%%%%%%%%%%%%
%%%%%%%%%%%%%%%%%%%%%%%%%%% Abstract  %%%%%%%%%%%%%%%%%%%%%%%%%%%%%%%%%%%%%%%%%%%%%%%%%%%
%%%%%%%%%%%%%%%%%%%%%%%%%%%%%%%%%%%%%%%%%%%%%%%%%%%%%%%%%%%%%%%%%%%%%%%%%%%%%%%%%%%%%%%%%%%%%%%%
\textbf {\large{Abstrakt}}\\

H\"{a}r skriver du en kort text om projektet\\\\






\textbf{Nyckelord:}

%	\include{Sections/Acknowledgements}
	\tableofcontents
	\newpage
%	\listoffigures
%	\newpage
%	\listoftables
%	\newpage
	
	%Here the actual report starts:
	%---------------------------------------------------------------------------------------------------------------
	\pagenumbering{arabic}
%	\include{Sections/Introduction}
%	\include{Sections/Technologies}
%	\include{Sections/CurrentNetwork}
%	\include{Sections/Achieving}
%	\include{Sections/ChosenTechnologies}
%	\include{Sections/TestEnvironment}
%	\include{Sections/TestResults}
%	\include{Sections/Discussion}
%	\include{Sections/Conclusions}

%
%	Ovanstående indelning med kapitel lagrade i en katalog har ersatts med allt i ett dokument.
%
	\Section{Introduktion}
		% \emph inneb�r emphasize, d.v.s. betona eller framh�ll -> kursiv stil
		\emph{Inled med lite text som beskriver innehållet i kapitlet.}
här kommer också övriga tecken

		% l�gg in en underrubrik (\subsection -> spillutrymme)
		\Subsection{Inledning/Bakrund}

		\Subsection{Tidigare forskning}

		\Subsection{Problemformulering}

		\Subsection{Syfte och frågeställning/hypotes}

		\Subsection{Avgränsning/Begränsning}

		\Subsection{ Målgrupp}

		\Subsection{Disposition}

\newpage

	\Section{Bakgrund/Teori}
		\emph{Inled med lite text som beskriver inneh\aa llet i kapitlet.}
		
		\Subsection{Rubrik 1}
			
		\Subsection{Rubrik 2}

		\Subsection{Rubrik etc}

\newpage			

	\Section{Metod}
		\emph{Inled med lite text som beskriver inneh\aa llet i kapitlet.}

		\Subsection{Vetenskaplig ansats}

		\Subsection{Datainsamling}

			\Subsubsection{Urval}

			\Subsubsection {Genomförande}
			
		\Subsection{Analys}

		\Subsection{Tillförlitlighet}

		\Subsection{Etiska överväganden}

\newpage

	\Section{Resultat/Emperi}
		\emph{Inled med lite text som beskriver inneh\aa llet i kapitlet.}
		
		\Subsection{Test 1}
			
		\Subsection{Test 2}	
\newpage
	
	\Section{Analys}
		\emph{Inled med lite text som beskriver inneh\aa llet i kapitlet.}
	
\newpage
	
	\Section{Diskussion}
		\emph{Inled med lite text som beskriver inneh\aa llet i kapitlet.}

		\Subsection{Problemlösning/resultat}

		\Subsection{Metodreflektion}

\newpage
		
	\Section{Avslutning}
		\emph{Inled med lite text som beskriver inneh\aa llet i kapitlet.}
		
		\Subsection{Slutsats}

		\Subsection{Förslag till fortsatt forskning}
\newpage

	
	%Here the bibliography starts:
	%---------------------------------------------------------------------------------------------------------------
	\newpage
	\clearpage
	\bibliographystyle{agsm}
	\bibliography{bachelor}
	\addcontentsline{toc}{section}{Referenser}
	\newpage
	
	%Here the end of the report starts - abbrivations, definitions and the last page:
	%---------------------------------------------------------------------------------------------------------------
%	\include{Sections/Definitions}
%	\addcontentsline{toc}{section}{Definitions}
%	\include{Sections/Abbreviations}
%	\addcontentsline{toc}{section}{Abbreviations}
%	\thispagestyle{empty}
\null
\vfill
\begin{table}[bp]
	\begin{adjustwidth}{-1in}{-1in}% adjust the L and R margins by 1 inch
		\begin{tabular}{p{4cm}l}
%		\includegraphics[scale=1.0]{img/lnuSymbol} & \includegraphics[scale=1.0]{img/lnuWordmark} \\
		& \\
		& 351 95 V\"{a}xj\"{o} / 391 82 Kalmar \\
		& Tel 0772-28 80 00 \\
		& dfm@lnu.se \\
		& Lnu.se/dfm
		\end{tabular}
	\end{adjustwidth}
\end{table}

%	\thispagestyle{empty}
\null
\vfill
\begin{table}[bp]
	\begin{adjustwidth}{-1in}{-1in}% adjust the L and R margins by 1 inch
		\begin{tabular}{p{4cm}l}
%		\includegraphics[scale=1.0]{img/lnuSymbol} & \includegraphics[scale=1.0]{img/lnuWordmark} \\
		& \\
		& 351 95 V\"{a}xj\"{o} / 391 82 Kalmar \\
		& Tel 0772-28 80 00 \\
		& dfm@lnu.se \\
		& Lnu.se/dfm
		\end{tabular}
	\end{adjustwidth}
\end{table}


\end{document}